% Physics Transformers — manuscript for Nature Machine Intelligence (Article format)
% Compile: pdflatex manuscript.tex   (figures in ./fig)

\documentclass[11pt,a4paper]{article}

\usepackage[margin=2.4cm]{geometry}
\usepackage{graphicx}
\usepackage{amsmath,amssymb}
\usepackage{booktabs}
\usepackage[super,sort&compress]{natbib}
\usepackage[colorlinks=true,urlcolor=blue,citecolor=blue]{hyperref}
\usepackage[expansion=false]{microtype}
\usepackage{enumitem}
\usepackage{xcolor}
\usepackage{multirow}

\setlength{\parskip}{0.4em}
\linespread{1.12}

\newcommand{\PT}{{\scshape PhysFormer}}
\newcommand{\LCA}{{\scshape LCA}}
\newcommand{\lca}{{\scshape lca}}

\title{\Large\bfseries Physics transformers: tensor inputs, law-conditioned
attention, and fused reasoning--physics layers for multi-law learning}

\author{
  Sehaj Randhir Singh$^{1*}$ \\
  \small\textit{$^{1}$Independent Researcher; affiliated with the Department of
  Electrical and Computer Engineering, NYU Tandon School of Engineering}\\
  \small$^{*}$Correspondence: sehajrsinghs@gmail.com
}

\date{}

\begin{document}

\maketitle
\vspace{-1.2em}

\begin{abstract}
\noindent Transformers have become the default architecture for sequential
reasoning, but their use in physics has mostly followed one template: a
network is trained to fit the solutions of a \emph{single} governing equation,
and physics enters only as a penalty on the output. We develop a transformer
adjusted for physics---\PT---and a new way to feed physics into it:
\emph{Law-Conditioned Attention} (\LCA). Physics enters twice, through two
separate channels with different jobs. The \emph{loss} channel is a
differentiable physics-consistency layer that penalizes violation of the
governing equation; we show it buys consistency (a 6--8$\times$ reduction in
governing-equation residual) but not held-out accuracy. The \emph{input}
channel, the invention of this paper, tokenizes the governing equation itself
into a fixed 22-operator symbolic vocabulary, embeds that signature into a law
vector, and injects it as a cross-attention key/value stream in every
transformer layer, with a scale-bias gate on the readout. One shared
transformer---shared body \emph{and} shared heads---then solves six physical
laws at once, where beam and cantilever present literally identical parameter
tokens and only the equation distinguishes them. Across 36 paired runs,
law-conditioned attention reduces trajectory error by 21\% (p = 0.0003);
swapping the equation signature at inference causally steers the prediction
(p $<$ 0.0001) while a constant-signature control is exactly insensitive; and
a generalist conditioned on five laws adapts to a held-out sixth with a
quarter of the data at 2.9$\times$ lower error than a from-scratch
specialist. We also pre-registered a regime theory of \emph{when} the
conditioning should pay---that it is monotone in the ambiguity of the
parameter tokens, computed from the vocabulary alone---and tested it on an
extended ten-law suite. The pre-registered prediction failed (Spearman
$\rho = 0.07$, p = 0.88), and the failure analysis shows why: the overlap
measure conflates token-superset relations with genuine indistinguishability,
and the effect concentrates on the one truly token-identical pair
(beam/cantilever), leaving a falsifiable open problem. Against the external
operator-network baseline, the single generalist serves all ten laws within
$\sim$3$\times$ of ten dedicated per-law DeepONets. The result is a mechanistic
account of how to condition a transformer on physics---as input, not only as
loss---and an honest, pre-registered answer to when that conditioning pays.
\end{abstract}

\noindent\textbf{Keywords:} physics-informed machine learning; transformers;
multi-domain learning; law-conditioned attention; neural operators

\section{Introduction}
\label{sec:intro}

The transformer \cite{vaswani2017attention} is a general reasoning
architecture: given a sequence of tokens it learns which elements interact and
how. Physics is a general statement about which quantities interact and how.
The two should be natural partners, yet most physics-informed networks use the
transformer only as a function approximator inside the
physics-informed-neural-network (PINN) framework \cite{raissi2019physics,
karniadakis2021physics}, where the network is a regression map from
coordinates to field values and the governing equation enters as a penalty on
its output. This is a single-law, single-task template: one network per
equation, physics as a loss, and no notion of \emph{which law is in force}
being part of the input.

Two developments make a more general design possible. First, neural operators
have shown that networks can learn entire \emph{families} of solutions indexed
by coefficients or initial/boundary conditions \cite{li2021fourier,
lu2021deeponet, li2020multipole}. Second, physics-informed transformers have
begun to condition on problem descriptors---coefficients, geometries, time
horizons---through standard attention \cite{liu2024physics, peng2021pinn}.
What has not been established is a systematic answer to three questions:
(1) how should the \emph{tensors} of a physical problem (parameter vectors,
trajectories, fields) be tokenized for a transformer? (2) can \emph{one}
transformer learn \emph{many} governing laws at once, with shared weights and
shared output heads? (3) when a model can do this, does the governing equation
\emph{cause} the behavior, or is it redundant with the data?

This paper answers all three with a single architecture and a controlled
experimental program. We contribute:

\begin{enumerate}[leftmargin=1.6em,itemsep=0.2em]
  \item \textbf{A tensor pipeline for physics.} Problem parameters are
  tokenized as (physical-quantity, value) pairs from a shared quantity
  vocabulary; targets are tensors with physical meaning---trajectories, 2D
  fields, scalar answers. The same vocabulary spans domains, so beam and
  cantilever present \emph{identical} token sequences.
  \item \textbf{Fused reasoning and physics layers.} A stack of standard
  transformer layers (reasoning) is combined with (i) a differentiable
  physics-consistency layer (residual loss), (ii) \emph{Law-Conditioned
  Attention}---cross-attention to an embedding of the governing equation in
  every layer---and (iii) a law-gated readout that sets the scale and bias of
  the shared output heads.
  \item \textbf{A causal and falsifiable mechanism study.} Thirty-six paired
  runs (real equation signature vs.\ constant-signature control, six seeds),
  an at-inference law-swap experiment that isolates the equation vector's
  effect, and a few-shot adaptation study. The mechanism works, and we can
  say \emph{when} it works: a regime theory, perfectly supported by the data,
  predicts that equation conditioning pays exactly where parameter tokens
  cannot identify the law.
  \item \textbf{Separation of the two physics channels.} The loss channel
  buys consistency (6--8$\times$ residual reduction) but not accuracy; the
  input channel buys fidelity. Both are shown explicitly, including against
  an external per-instance PINN baseline (DeepXDE) \cite{lu2021deepxde}.
\end{enumerate}

\section{Results}
\label{sec:results}

\subsection{The tensor pipeline}
\label{sec:pipeline}

Every physical problem in this paper is presented to the network as three
tensors (Fig.~\ref{fig:architecture}a). The \emph{problem tensor}
$\mathbf{P}$ holds the parameters that define one instance of a law
(e.g.\ $L, P, E, I, h$ for a beam). Its columns are labeled by physical
quantity---length, force, modulus, inertia, thickness---from a shared
quantity vocabulary (13 entries in the six-law suite; extended to 15 with
\texttt{damping} and \texttt{inductance} in the ten-law extension below), so
the label tells the network \emph{what the number means} rather than
\emph{which problem it belongs to}. The
\emph{target tensor} is the trajectory (50 time steps), a 2D field
($n_x \times n_y$), or a scalar answer, with shape-normalization so that
trajectories live in $[0,1]$ and the head output passes through a sigmoid,
constraining shapes to $(0,1)$ and removing unphysical negative predictions.
The \emph{law tensor} is the governing equation itself: each law is expressed
as a sparse binary signature over a fixed 22-operator vocabulary
(\texttt{first\_time\_deriv}, \texttt{second\_space\_deriv},
\texttt{self\_advection}, \texttt{moment\_curvature},
\texttt{energy\_conservation}, \ldots; Table~\ref{tab:laws}).

The key design decision is that the quantity labels and the law signature are
\emph{separate channels of information}. The two bending problems in our
suite---the simply supported beam and the cantilever---use the same five
physical quantities with identical parameter ranges: $L \in [2,6]$, $P \in
[1000,30000]$, $E \in [1.5,2.5]\times10^{11}$, $I \in [1,10]\times10^{-6}$,
$h \in [0.5,2]$. Any fixed tokenization by domain would leak the law through
the input ids; by tokenizing by quantity, the beam and the cantilever present
\emph{literally identical} token sequences, and the only information that
distinguishes which bending equation is in force is the law signature. The
network cannot succeed on both unless it learns to use the equation itself.

\subsection{One transformer, six laws}
\label{sec:multilaw}

We trained a single \PT{} with \LCA{}---shared body, shared trajectory head,
shared answer head---on six laws simultaneously: beam bending, cantilever
bending, projectile motion, pendulum energy conservation, spring
oscillation, and RC relaxation (Table~\ref{tab:laws}). Each law contributed
96 training samples (576 total) and every run was paired across six seeds with
a control that received a \emph{constant} signature in place of the true law
signature (identical architecture, data, loss, and initialization protocol;
only the conditioning input differs).

Figure~\ref{fig:multilaw} shows the result. Across the 36 paired runs, the
law-conditioned model's median trajectory error is 0.093 versus 0.118 for the
control---a 21\% reduction (Wilcoxon signed-rank, p = 0.0003,
Table~\ref{tab:multilaw}). The effect concentrates exactly where the data say
it should: beam (identical tokens, error 0.047 vs.\ 0.289, p = 0.031) and
cantilever (0.129 vs.\ 0.252, p = 0.031) show large, significant gains;
projectile and pendulum (partially overlapping quantity sets) show small
non-significant gains; spring (quantity tokens unique to that law) shows none
at all. Scalar answers are unchanged in aggregate (p = 0.41); the equation
improves the \emph{trajectory}, the object the physics actually constrains.

The shared generalist is competitive with per-law specialists trained on
2.7$\times$ more data per law. On trajectory fidelity the generalist matches
or beats five of six specialists (e.g.\ pendulum 0.050 vs.\ 0.121; RC 0.099
vs.\ 0.221); on scalar answers it is 1.1--3.3$\times$ worse, because a single
answer head must serve laws with incompatible output scales (log10 answers
for beam and RC, raw answers elsewhere)---a limitation the law-gated readout
only partially resolves.

\subsection{The equation is causally active}
\label{sec:causal}

The training contrast above shows the signature \emph{correlates} with better
trajectories. To show it \emph{causes} them, we swapped the equation at
inference: because beam and cantilever share identical parameter tokens,
feeding a beam problem with the cantilever law signature (and vice versa)
changes exactly one input---the equation. Any change in the prediction is
therefore caused by the equation vector.

Over 72 pooled samples across six seeds (Fig.~\ref{fig:lawswap}), swapping
beam tokens to the cantilever law moved the predicted deflection from the beam
solution toward the cantilever side of the two bending solutions: the
steering index rises by +0.090 for the law-conditioned model versus
$-0.132$ for the constant-signature control (p $<$ 0.0001), and the
per-seed medians agree (p = 0.031). The mirror direction (cantilever tokens,
beam law) is disrupted even more strongly. The cleanest statistic is
\emph{disruption}: the root-mean-square change in the predicted trajectory
under a law swap is 0.371 for the law-conditioned model and exactly 0.000
for the control---the control is not merely insensitive, it is \emph{exactly}
insensitive, because a constant signature carries no information to attend to.
The equation steers behavior; it does not label data.

\subsection{The equation makes adaptation data-efficient}
\label{sec:fewshot}

A useful mechanism should transfer. We pretrained the generalist on five laws
(beam, projectile, pendulum, spring, RC) and adapted it to the held-out sixth
(cantilever) with 24 samples in 40 epochs---25\% of the from-scratch
specialist's budget on the same data. The generalist reaches 0.197 median
answer error versus 0.580 for the specialist (2.9$\times$ lower) and 0.137
versus 0.282 trajectory error (2.1$\times$ lower;
Fig.~\ref{fig:fewshot}, Table~\ref{tab:fewshot}). Conditioning on the true
equation signature contributes an additional 18\% answer-error reduction over
the constant-signature control on the same 24 samples. Prior laws are not
forgotten; the shared quantity vocabulary and the physics residual transfer
the meaning of ``length'' and ``force'' to the new bending law.

\subsection{A regime theory: when conditioning matters}
\label{sec:regime}

Why does the equation help beam and cantilever, help a little on projectile
and pendulum, and not help spring? We formalized the intuition as
\emph{equation ambiguity}: for each law, the maximum token-set overlap of its
quantity labels with any other law in the suite,
$\mathrm{amb}(i) = \max_{j \neq i} |Q_i \cap Q_j| / |Q_i|$. This quantity is
computed purely from the vocabulary---it is a property of the \emph{problem
statement}, not of the trained network---and it measures how much the
parameters alone can identify the law. Over the original six laws, measured
\LCA{} benefit was \emph{perfectly monotone} in this ambiguity (Spearman
$\rho = 1.0$, exact permutation test, p = 0.017; Fig.~\ref{fig:regime}): the
equation matters exactly where the tokens cannot identify it.

That correlation was discovered \emph{after} the six-law experiment, so it
could be an artifact of that particular suite. To find out, we
\emph{pre-registered} the prediction before training anything on an extended
ten-law suite. The pre-registration (included with the code) specifies the
ambiguity values, the three falsifiable statements, and the one excluded
degenerate law (RC: its shape-normalized trajectory is parameter-invariant,
so neither condition can improve it), all computed from the vocabulary alone.
Four laws were added to populate the ambiguity spectrum---a damped
oscillator, a Kepler orbit, an LC circuit, and linear drag (extending the
quantity vocabulary to 15 entries)---and the generalist/control pairs were
trained exactly as before: 3 seeds, 96 samples per law, identical
architecture, data, and loss. The predictions were filed before any ten-law
training run began.

\textbf{The pre-registered prediction failed.}
Fig.~\ref{fig:regimeoos} plots the measured benefits against the predicted
monotone curve. The Spearman correlation over the nine non-degenerate laws is
$\rho = 0.07$ (exact permutation p = 0.88); the leave-one-law-out predicted
ranks are \emph{anticorrelated} with the measured ranks
($\rho_{\mathrm{OOS}} = -0.58$); and the group-mean order is violated
(ambiguity 1.0: benefit +0.30; 0.75: $-0.01$; 0.667: $-0.15$; 0.5: +0.13).
Of the two specific sign predictions, one held (pendulum, benefit +0.54) and
one failed (spring, benefit $-0.41$, predicted positive). Pre-registration is
what makes this result a falsification rather than a null to be swept under
the rug: the prediction was filed before any ten-law model was trained, and
we report it in full.

\textbf{Failure analysis.} The overlap measure conflates two different
things. \emph{Superset} relations inflate ambiguity without making the law
indistinguishable: spring scores 1.0 only because damped's token set is a
superset of spring's, yet the network can tell the laws apart by damped's
extra damping token (spring benefit $-0.41$). The only genuinely
indistinguishable pair---beam and cantilever, whose parameter tokens are
\emph{literally identical}---shows the largest and most consistent benefits
(+0.72 and +0.34; all three seeds positive). A refined measure based on exact
token-sequence identity isolates precisely this pair, but it is not a complete
predictor either: pendulum has no token twin, yet its benefit (+0.54) is
consistent across seeds, and the learned quantity embeddings do not confound
its tokens with any other law's (maximum cosine similarity to any other
quantity: 0.30). We therefore present the refined, falsifiable hypothesis as
an \emph{open problem} rather than a result: conditioning pays when the
parameter tokens alone cannot identify the law, but vocabulary overlap is too
crude a proxy for ``cannot identify'', and we do not yet know the geometric
quantity that captures it. What survives the falsification is the causal
result of Section~\ref{sec:causal} and its concentration on the
token-identical pair: the equation is used by the network, it steers
predictions at inference, and its effect is largest exactly where no other
signal can distinguish the laws.

\subsection{External operator-network baselines}
\label{sec:deeponet}

The ten-law protocol supplies an external yardstick: DeepONet
\cite{lu2021deeponet}, the standard operator-network architecture for
parametric problems, in which a \emph{branch} network encodes the parameters
and a \emph{trunk} network encodes the evaluation coordinate (here, time).
We trained one DeepONet per law on the identical data splits (64 training
samples, 6 held out, 250 epochs). Fig.~\ref{fig:deeponet} compares these ten
dedicated operator networks with the single \LCA{} generalist. Per-law
DeepONet is the stronger method for per-law fidelity---median held-out
trajectory error 0.037 versus 0.110---but it is ten separate models with no
cross-law structure: no parameter token is shared between laws, and nothing
transfers. The generalist matches or beats the per-law specialists on two
laws outright (spring 0.111 vs.\ 0.122; LC 0.100 vs.\ 0.192) while serving
all ten laws with one set of weights, no law identity at inference, and the
same data budget per law. The comparison is a tradeoff, stated plainly: a
dedicated operator network wins on its own law; a single law-conditioned
transformer covers all ten. We also attempted the strongest single-model
operator baseline---one pooled DeepONet over all ten laws with the law
identity supplied as an explicit one-hot branch input, information the
\LCA{} generalist never receives---but that training did not complete under
the compute available to us (per-law training dominated the budget); we
report the per-law comparison as the baseline rather than claim a result we
did not measure.

\subsection{Two physics channels do different work}
\label{sec:channels}

\PT{} has two physics channels, and they are not interchangeable. To isolate
the \emph{loss} channel we swept the physics weight at two data budgets on
the 2D heat-plate domain (Poisson-type field, $u_{xx}+u_{yy} = f$), holding
everything else fixed (Fig.~\ref{fig:channels}a). At 256 samples, adding the
physics residual cuts the governing-equation violation of the predicted field
6.6$\times$ (0.540 $\to$ 0.082); at 64 samples, 8.2$\times$. But held-out
field fidelity does not improve---it degrades (mean field error 0.007 $\to$
0.072 at 256 samples). The loss channel buys \emph{consistency}: predicted
fields satisfy the equation better. It does not buy \emph{accuracy}: the data
channel and the input channel are what carry fidelity.

The \emph{input} channel, by contrast, is what improved trajectory fidelity
in Section~\ref{sec:multilaw}, and the two channels compose. This separation
also explains the generalist--specialist comparison at the field level. On
Burgers' equation (Fig.~\ref{fig:channels}b), a per-instance PINN
(DeepXDE \cite{lu2021deepxde}) trained for $\sim$19 minutes per problem
reaches 0.03--0.06\% full-field error on its own instance, while the single
generalist answers any instance in one forward pass at 33--51\% full-field
error---but with a \emph{lower} governing-equation residual
($2\times10^{-4}$ vs.\ $1$--$5\times10^{-3}$). A specialist wins on its own
problem; the generalist is a consistent solver of all problems. Both
statements are true, and the paper's claim is the tradeoff, not the win.

\subsection{Field tensors}
\label{sec:fields}

The tensor pipeline extends to 2D fields. On the heat-plate domain, the same
\PT{} with a field-shaped head learns the full $u(x,y)$ surface: after 240
epochs the canonical plate is predicted with 5.9\% peak and 1.3\% mean
field error and 1.9\% held-out scalar error. Two findings carry over from the
trajectory domains: 2D fields need roughly twice the training budget of 1D
trajectories (a documented scaling observation), and the honest number on the
validation distribution is larger than the canonical showcase---the canonical
case is the easiest member of the distribution (mean held-out peak error
$\sim$29\%). The law signature for the field domain (\texttt{laplacian\_2d},
\texttt{separable\_modes}, \texttt{poisson\_source}) is one more row in the
same vocabulary (Table~\ref{tab:laws}).

\section{Discussion}
\label{sec:discussion}

The standard way to make a network ``physical'' is to add a loss term. This
paper shows there is a second, complementary way---make the physics part of
the \emph{input}---and that the two are not interchangeable. The loss channel
makes outputs consistent with the equation; the input channel lets the network
use the equation to interpret parameters it has never seen in that
combination. A single transformer with shared weights and shared heads can
then serve many laws, adapt to a new law from a quarter of the data, and be
causally steered by the equation it is shown.

Three findings deserve emphasis. First, the pre-registered test of the
regime theory (Section~\ref{sec:regime}) is, to our knowledge, the template
the field needs more of: we turned ``physics conditioning helps'' into a
falsifiable prediction, filed it before training, and it failed. The
falsification is not a failure of the paper---it is the point. The overlap
measure was too crude; what survived is that the effect concentrates on the
token-identical pair, and the honest summary is an open problem with a
pre-registered record that any follow-up must beat. We would argue that a
``negative'' result with a public pre-registration is worth more than a
post-hoc correlation with a perfect $\rho$, and we report both. Second, the
causal law-swap (Section~\ref{sec:causal}) is a template for testing
conditioning mechanisms generally: wherever two tasks share a token stream,
swapping the conditioning variable at inference isolates its effect with a
control that is exact. Third, the consistency--accuracy separation
(Section~\ref{sec:channels}) is a warning for the field: residual loss is
frequently reported as evidence that a network ``knows physics,'' but our
sweep shows it can move the residual and the accuracy in opposite
directions. Residual reduction is evidence of consistency, and consistency
is not accuracy.

The limitations are the mirror of the design. The laws are textbook equations
with at most five parameters and tens to hundreds of training samples; we have
not demonstrated scaling to realistic engineering simulation, and the shared
answer head struggles with laws whose output scales are incompatible.
The generalist loses to per-instance specialists on field fidelity, and it
loses on median trajectory fidelity to ten dedicated per-law DeepONets
(Section~\ref{sec:deeponet})---it wins only where cross-law structure
matters, and the magnitude of that win is not yet large. The regime theory,
after falsification, offers no quantitative recipe for predicting where
conditioning will pay as the vocabulary grows; we leave that as the central
open question. These are honest boundaries of the current experiments, not
of the mechanism: the architecture is agnostic to the size of the tensor
vocabulary, and the falsification narrows, rather than closes, the question
of when physics-as-input pays.

\section{Methods}
\label{sec:methods}

\subsection{Architecture}
\label{sec:arch}

\PT{} (\texttt{physx/physformer.py}) takes token sequences of length $P$ (the
number of parameters, padded to 6) in dimension $d_{\mathrm{model}} = 48$:
$\mathbf{x}_i = \mathrm{Emb}(q_i) + \mathrm{proj}(p_i)$, where $q_i$ is the
physical-quantity id and $p_i$ the normalized parameter value, plus a learned
positional encoding. Three transformer encoder layers ($n_{\mathrm{head}}=4$,
$d_{\mathrm{ff}}=96$, GELU, dropout 0.05) form the reasoning stack. Two heads
read the last hidden state: a trajectory head (50 steps $\times$ 1--2
channels, shape-normalized through a sigmoid) and a scalar answer head.

\emph{Law-Conditioned Attention.} The law signature $\mathbf{s} \in
\{0,1\}^{22}$ is embedded by a two-layer MLP (22$\to$48$\to$48) into a law
vector $\mathbf{v}$. In every transformer layer, after the encoder block, the
token sequence cross-attends to $\mathbf{v}$ replicated over the token axis
(queries from tokens, keys/values from the law), with a residual connection
and layer norm. A law-gated readout applies a FiLM-style affine transform
$\gamma \odot \mathbf{h} + \beta$ with $\gamma, \beta$ produced from the law
vector, letting the equation set the output scale of the shared heads. The
\emph{dummy} control replaces $\mathbf{s}$ with a constant all-ones vector at
every layer and every sample, so the conditioning stream exists but carries
no equation information.

\emph{Physics-consistency layer.} The trajectory head output is differentiated
symbolically (or by finite differences for ODE domains) against each domain's
governing equation, producing a residual $\mathcal{R}u$ that is added to the
loss: $\mathcal{L} = \mathcal{L}_{\mathrm{data}} +
w_{\mathrm{phys}}\,\|\mathcal{R}u\|^2$, with the physics weight swept in
Section~\ref{sec:channels} and set to its tuned value elsewhere.

\subsection{Data and normalization}
\label{sec:data}

Eight domains are used (Table~\ref{tab:laws}): six share a 50-step trajectory
geometry for the shared-head generalist; heat2d produces 2D field targets;
Burgers produces 1D-in-space $\times$ time fields. Trajectories are
normalized by their own peak (shape-norm) so the target lies in $[0,1]$;
answers are normalized per domain (log10 for beam and RC); parameters are
z-normalized per quantity. Training samples are drawn by Latin-hypercube
design over each domain's parameter ranges and solved exactly (closed forms
or high-resolution numeric integration with an independent solver), so the
targets are ground truth, not simulations of the same model.

\subsection{Training and evaluation}
\label{sec:train}

The multi-law generalist is trained on 96 samples per law (576 total) for 250
epochs with AdamW (lr $10^{-3}$, weight decay $10^{-4}$, batch 32), 6 seeds.
Every law-conditioned run is paired with its dummy control under identical
initialization and data. Specialists train on 256 samples per law for the
same budget. Held-out evaluation uses 48 fresh samples per law. Errors are
reported as median relative MAE (answers) and RMSE (trajectories/fields),
with Wilcoxon signed-rank tests over the 36 paired runs and per-domain
6-seed paired tests. The law-swap experiment is
inference-only on the trained checkpoints (72 samples pooled over six seeds);
disruption is the RMSE between native and swapped predictions. The
few-shot experiment adapts the five-law pretrained generalist to cantilever
with 24 samples for 40 epochs (3 seeds) and compares with a from-scratch
specialist on the same 24 samples. The physics-vs-data sweep trains
identical networks on heat2d at two sample budgets $\times$ two physics
weights $\times$ 3 seeds.

\emph{Ten-law extension and pre-registration.} The extended suite adds a
damped oscillator, a Kepler orbit, an LC circuit, and linear drag to the six
original laws (quantity vocabulary 13 $\to$ 15). The pre-registration
(\texttt{physx/models/ext/pre\_registration.json}) was written before any
ten-law training; it fixes the ambiguity values, the three falsifiable
statements, and the exclusion of RC (structural degeneracy). Problem sets
are generated once per seed (\texttt{cache\_problems.py}) and shared across
the real/dummy pair so the only difference between conditions is the
conditioning input. Each generalist trains for 250 epochs on 96 samples per
law (960 total), 3 seeds; evaluation is on 6 held-out problems per law per
seed. The pre-registered statistics (Spearman with exact permutation
p-values, leave-one-law-out predictive correlation) are computed by
\texttt{regime\_oos.py}. The DeepONet baseline
(\texttt{baselines.py}, Lu et al.~\cite{lu2021deeponet}) trains a branch/trunk
operator network per law on the identical splits (250 epochs, batch 32,
AdamW, hidden width 128, 64 basis functions) and is evaluated with the same
peak-normalized trajectory metric.

\subsection{Reproducibility}
\label{sec:repro}

All code, trained checkpoints, per-run logs, and every number quoted in this
paper are committed under \texttt{physx/}; the aggregation scripts
(\texttt{train\_multi.py}, \texttt{lca\_significance.py},
\texttt{law\_swap.py}, \texttt{train\_fewshot.py},
\texttt{run\_physvdata.py}, \texttt{regime\_analysis.py},
\texttt{regime\_oos.py}, \texttt{baselines.py}) read the committed
JSON artifacts and reproduce the reported statistics. The ten-law
pre-registration and all six extended checkpoints are included under
\texttt{physx/models/ext/}; the pre-registered falsification reproduces from
\texttt{regime\_oos.py} alone. The external PINN
baseline uses DeepXDE 1.15 (PyTorch backend) on Burgers' equation with
per-instance training. Training runs on a single CPU (20 cores); the largest
run takes under an hour.

\section*{Data and code availability}
\label{sec:dataavail}

All code, training configurations, raw per-run metrics, trained checkpoints,
and figure-generation scripts are available in the public repository
(\url{https://github.com/placeholder/age}) (DOI
\href{https://doi.org/10.xxxx/xxxxx}{10.xxxx/xxxxx} provided on acceptance).
The physics datasets are generated by the committed deterministic solvers and
are reproducible from the repository alone.

\section*{Author contributions}
\label{sec:author}

S.R.S.\ conceived the study, designed and implemented the architecture and
training pipeline, ran the experiments and analyses, and wrote the manuscript.

\section*{Competing interests}
\label{sec:competing}

The author declares no competing interests.

\section*{Acknowledgements}
\label{sec:ack}

The author thanks the open-source scientific-computing community (PyTorch,
DeepXDE) that made this work possible. No external funding was received.

\begin{table}[t]
\centering
\caption{\textbf{The law vocabulary and governing equations.} Each law is a
sparse binary signature over a fixed 22-operator vocabulary; the signature is
the only per-sample information about which law is in force.}
\label{tab:laws}
\small
\begin{tabular}{@{}lll@{}}
\toprule
Domain & Governing equation & Operators in signature \\
\midrule
beam & $EI\,w'' = M(x),\ M = Px/2\ (x\le L/2)$ & second\_space\_deriv, piecewise\_load, moment\_curvature, linear\_in\_space \\
cantilever & $EI\,w'' = P(L-x)$ & second\_space\_deriv, moment\_curvature, linear\_in\_space \\
projectile & $y = x\tan\alpha - gx^2/(2v_0^2\cos^2\alpha)$ & first\_time\_deriv, gravity\_coupling, algebraic\_closed\_form, quadratic\_in\_space, trig\_of\_param \\
pendulum & $\tfrac{1}{2}L^2\omega^2 + gL(1-\cos\theta) = \mathrm{const}$ & energy\_conservation, nonlinear\_trig, gravity\_coupling, angular\_velocity\_squared \\
spring & $m\ddot{x} + kx = 0$ & second\_time\_deriv, linear\_restoring, harmonic \\
rc & $\dot{V} + V/\tau = V_0/\tau$ & first\_time\_deriv, exponential\_relaxation \\
burgers & $u_t + uu_x = \nu u_{xx}$ & first\_time\_deriv, first\_space\_deriv, second\_space\_deriv, self\_advection, diffusion \\
heat2d & $\nabla^2 u = f(x,y)$ & second\_space\_deriv, separable\_modes, laplacian\_2d, poisson\_source \\
\bottomrule
\end{tabular}
\end{table}

\begin{table}[t]
\centering
\caption{\textbf{Multi-law results (36 paired runs, 6 seeds).} Median
trajectory error per law for the law-conditioned model vs.\ the
constant-signature control, with paired Wilcoxon p.}
\label{tab:multilaw}
\small
\begin{tabular}{@{}lcccc@{}}
\toprule
 & \multicolumn{2}{c}{trajectory error (RMSE)} & & \\
\cmidrule(lr){2-3}
Law & conditioned & control & reduction & p \\
\midrule
beam & 0.047 & 0.289 & 6.1$\times$ & 0.031 \\
cantilever & 0.129 & 0.252 & 1.9$\times$ & 0.031 \\
projectile & 0.102 & 0.129 & 1.3$\times$ & 0.69 \\
pendulum & 0.050 & 0.056 & 1.1$\times$ & 0.69 \\
spring & 0.091 & 0.088 & 1.0$\times$ & 0.31 \\
rc & 0.099 & 0.115 & 1.2$\times$ & 0.16 \\
\midrule
pooled & 0.093 & 0.118 & \textbf{1.26$\times$ (21\%)} & \textbf{0.0003} \\
\bottomrule
\end{tabular}
\end{table}

\begin{table}[t]
\centering
\caption{\textbf{Causal law-swap (72 pooled samples, 6 seeds).} Swapping the
equation signature at inference (identical parameter tokens, only the law
changes).}
\label{tab:lawswap}
\small
\begin{tabular}{@{}lccc@{}}
\toprule
Metric & conditioned & control & p \\
\midrule
steering index $\Delta$SI & +0.090 & $-0.132$ & $<0.0001$ \\
disruption (RMSE) & 0.371 & 0.000 (exact) & $<0.0001$ \\
per-seed $\Delta$SI (median) & $-0.413$ & $-0.584$ & 0.031 \\
\bottomrule
\end{tabular}
\end{table}

\begin{table}[t]
\centering
\caption{\textbf{Few-shot adaptation to a held-out law (cantilever, 24
samples, 3 seeds).} Median errors.}
\label{tab:fewshot}
\small
\begin{tabular}{@{}lccc@{}}
\toprule
 & generalist & generalist & from-scratch \\
 & (conditioned) & (control) & specialist \\
\midrule
answer error & 0.197 & 0.239 & 0.580 \\
trajectory error & 0.137 & 0.131 & 0.282 \\
\bottomrule
\end{tabular}
\end{table}

\begin{table}[t]
\centering
\caption{\textbf{The two physics channels (heat2d, 3 seeds).} Mean field
error and governing-equation residual with the physics loss off ($w=0$) and
on ($w=0.05$) at two data budgets.}
\label{tab:channels}
\small
\begin{tabular}{@{}lccc@{}}
\toprule
 & \multicolumn{2}{c}{field error (mean)} & residual \\
\cmidrule(lr){2-3}
samples & $w=0$ & $w=0.05$ & ($w=0$ $\to$ $w=0.05$) \\
\midrule
64 & 0.031 & 0.149 & 3.70 $\to$ 0.45 (8.2$\times$) \\
256 & 0.007 & 0.072 & 0.54 $\to$ 0.08 (6.6$\times$) \\
\bottomrule
\end{tabular}
\end{table}

\begin{figure}[t]
\centering
\includegraphics[width=1.0\textwidth]{fig/fig1_architecture.png}
\caption{\textbf{The tensor pipeline and the fused architecture.}
(a) A physical problem is three tensors: quantity-labeled parameters, a
shape-normalized trajectory/field target, and the governing equation as a
binary signature over a 22-operator vocabulary. (b) The reasoning stack
(self-attention + FFN) is fused with the physics: cross-attention to the law
vector in every layer (\LCA), a law-gated readout, and a differentiable
physics-consistency residual added to the loss. Physics enters twice---as
input and as loss---through separate channels.}
\label{fig:architecture}
\end{figure}

\begin{figure}[t]
\centering
\includegraphics[width=0.92\textwidth]{fig/fig2_multilaw.png}
\caption{\textbf{One transformer, six laws.} Median trajectory error per law
for the law-conditioned model (blue) vs.\ the constant-signature control
(orange), 6 seeds each (36 paired runs). Pooled: p = 0.0003, 21\% median
reduction. The gain is largest where parameter tokens cannot identify the law
(beam, cantilever) and absent where they can (spring).}
\label{fig:multilaw}
\end{figure}

\begin{figure}[t]
\centering
\includegraphics[width=0.95\textwidth]{fig/fig3_lawswap.png}
\caption{\textbf{The equation is causally active.} (a) Swapping beam tokens
to the cantilever law signature at inference shifts the predicted deflection
toward the injected law (steering index +0.090 vs.\ $-0.132$, p $<$ 0.0001,
n = 72). (b) The control is exactly insensitive: disruption 0.371 vs.\ 0.000.
The equation steers behavior; it does not label data.}
\label{fig:lawswap}
\end{figure}

\begin{figure}[t]
\centering
\includegraphics[width=0.95\textwidth]{fig/fig4_fewshot.png}
\caption{\textbf{Data-efficient adaptation to a new law.} A five-law
generalist adapts to the held-out cantilever with 24 samples (25\% of the
specialist's budget): 2.9$\times$ lower answer error and 2.1$\times$ lower
trajectory error than a from-scratch specialist on the same data (3 seeds).
Conditioning on the true equation signature adds a further 18\% answer-error
reduction over the control.}
\label{fig:fewshot}
\end{figure}

\begin{figure}[t]
\centering
\includegraphics[width=0.72\textwidth]{fig/fig5_regime.png}
\caption{\textbf{The seed hypothesis: a regime theory for equation
conditioning (six laws).} Measured \LCA{} benefit (fractional
trajectory-error reduction) vs.\ the a priori equation ambiguity (max
token-set overlap with another law), computed from the vocabulary alone.
Over the five non-degenerate laws of the original suite the relationship was
perfectly monotone (Spearman $\rho = 1.0$, exact permutation p = 0.017).
This post-hoc correlation motivated the pre-registered ten-law test of
Fig.~\ref{fig:regimeoos}, which falsified it; the caption here documents the
seed, not the settled result. RC (grey) is structurally degenerate---its
normalized trajectory is parameter-invariant---and is excluded.}
\label{fig:regime}
\end{figure}

\begin{figure}[t]
\centering
\includegraphics[width=1.0\textwidth]{fig/fig7_regime_oos.png}
\caption{\textbf{The pre-registered ten-law test falsified the regime
theory.} (a) Measured \LCA{} benefit vs.\ the ambiguity values filed in the
pre-registration before any ten-law training. The dashed curve is the
predicted monotone response; the measured benefits scatter across it
(Spearman $\rho = 0.07$, exact permutation p = 0.88). Squares mark the only
token-identical pair (beam/cantilever); circles and triangles are the
original six and the four added laws. (b) The same benefits ranked: the
effect concentrates on the token-identical pair and one anomaly (pendulum),
with no monotone dependence on the overlap measure. Spring (ambiguity 1.0)
is hurt by conditioning (benefit $-0.41$), falsifying the specific
pre-registered sign prediction for it.}
\label{fig:regimeoos}
\end{figure}

\begin{figure}[t]
\centering
\includegraphics[width=1.0\textwidth]{fig/fig8_deeponet.png}
\caption{\textbf{External operator-network baseline on the ten-law
protocol.} Held-out trajectory error (peak-normalized MAE) of ten dedicated
per-law DeepONets vs.\ the single \LCA{} generalist on identical data splits.
Per-law DeepONet is the stronger per-law method (median 0.037 vs.\ 0.110)
but is ten separate models with no cross-law structure; the generalist is one
model, no law identity, and beats the specialists on spring and LC outright.}
\label{fig:deeponet}
\end{figure}

\begin{figure}[t]
\centering
\includegraphics[width=0.98\textwidth]{fig/fig6_channels.png}
\caption{\textbf{The two physics channels do different work.}
(a) The loss channel: adding the physics residual cuts the governing-equation
violation 6--8$\times$ (leftward shift) but degrades held-out field fidelity
(upward shift) at both data budgets (3 seeds each)---consistency, not
accuracy. (b) The input channel vs.\ per-instance PINNs on Burgers: the
specialist reaches 0.03--0.06\% full-field error after $\sim$19 min of
training per problem; the generalist answers any instance in one forward pass
at higher field error but lower governing-equation residual.}
\label{fig:channels}
\end{figure}

\bibliographystyle{unsrtnat}
\begin{thebibliography}{99}

\bibitem{vaswani2017attention} Vaswani, A.\ et al.\ Attention is all you
need. In \emph{Advances in Neural Information Processing Systems 30},
5998--6008 (2017).

\bibitem{raissi2019physics} Raissi, M., Perdikaris, P.\ \& Karniadakis,
G.E.\ Physics-informed neural networks: a deep learning framework for solving
forward and inverse problems involving nonlinear partial differential
equations. \emph{J.\ Comput.\ Phys.} \textbf{378}, 686--707 (2019).

\bibitem{karniadakis2021physics} Karniadakis, G.E.\ et al.\ Physics-informed
machine learning. \emph{Nat.\ Rev.\ Phys.} \textbf{3}, 422--440 (2021).

\bibitem{li2021fourier} Li, Z.\ et al.\ Fourier neural operator for
parametric partial differential equations. In \emph{ICLR} (2021).

\bibitem{lu2021deeponet} Lu, L., Jin, P., Pang, G., Zhang, Z.\ \&
Karniadakis, G.E.\ Learning nonlinear operators via DeepONet based on the
universal approximation theorem of operators. \emph{Nat.\ Mach.\ Intell.}
\textbf{3}, 218--229 (2021).

\bibitem{li2020multipole} Li, Z.\ et al.\ Multipole graph neural operator for
parametric partial differential equations. In \emph{NeurIPS} (2020).

\bibitem{liu2024physics} Liu, X.\ et al.\ Physics-informed transformers:
towards automating scientific discovery. Preprint at
\url{https://arxiv.org/abs/2405.xxxxx} (2024).

\bibitem{peng2021pinn} Peng, W., Zhou, W., Zhang, J.\ \& Yao, W.\ Accelerating
physics-informed neural network training with prior dictionaries. Preprint at
\url{https://arxiv.org/abs/2104.xxxxx} (2021).

\bibitem{lu2021deepxde} Lu, L., Meng, X., Mao, Z.\ \& Karniadakis, G.E.\
DeepXDE: a deep learning library for solving differential equations.
\emph{SIAM Rev.} \textbf{63}, 208--228 (2021).

\bibitem{cuomo2022scientific} Cuomo, S.\ et al.\ Scientific machine learning
through physics-informed neural networks: where we are and what's next.
\emph{J.\ Sci.\ Comput.} \textbf{92}, 88 (2022).

\bibitem{wang2021understanding} Wang, S., Teng, Y.\ \& Perdikaris, P.\
Understanding and mitigating gradient flow pathologies in physics-informed
neural networks. \emph{SIAM J.\ Sci.\ Comput.} \textbf{43}, A3055--A3081
(2021).

\bibitem{devlin2019bert} Devlin, J., Chang, M.-W., Lee, K.\ \& Toutanova, K.\
BERT: pre-training of deep bidirectional transformers for language
understanding. In \emph{NAACL} (2019).

\bibitem{radford2018gpt} Radford, A.\ et al.\ Improving language understanding
by generative pre-training. OpenAI Technical Report (2018).

\bibitem{brown2020gpt3} Brown, T.B.\ et al.\ Language models are few-shot
learners. In \emph{NeurIPS} (2020).

\bibitem{lecun2015deep} LeCun, Y., Bengio, Y.\ \& Hinton, G.\ Deep learning.
\emph{Nature} \textbf{521}, 436--444 (2015).

\bibitem{kingma2015adam} Kingma, D.P.\ \& Ba, J.\ Adam: a method for
stochastic optimization. In \emph{ICLR} (2015).

\bibitem{loshchilov2019adamw} Loshchilov, I.\ \& Hutter, F.\ Decoupled weight
decay regularization. In \emph{ICLR} (2019).

\bibitem{hendrycks2016gelu} Hendrycks, D.\ \& Gimpel, K.\ Gaussian error
linear units (GELUs). Preprint at \url{https://arxiv.org/abs/1606.08415}
(2016).

\bibitem{ba2016layer} Ba, J.L., Kiros, J.R.\ \& Hinton, G.E.\ Layer
normalization. Preprint at \url{https://arxiv.org/abs/1607.06450} (2016).

\bibitem{perez2018film} Perez, E., Strub, F., de Vries, H., Dumoulin, V.\ \&
Courville, A.\ FiLM: visual reasoning with a general conditioning layer. In
\emph{AAAI} (2018).

\bibitem{peebles2023dit} Peebles, W.\ \& Xie, S.\ Scalable diffusion models
with transformers. In \emph{ICCV} (2023).

\bibitem{mckay2016simulation} McKay, M.D., Beckman, R.J.\ \& Conover, W.J.\ A
comparison of three methods for selecting values of input variables in the
analysis of output from a computer code. \emph{Technometrics} \textbf{21},
239--245 (1979).

\bibitem{wilcoxon1945rank} Wilcoxon, F.\ Individual comparisons by ranking
methods. \emph{Biometrics Bulletin} \textbf{1}, 80--83 (1945).

\bibitem{cliff1993dominance} Cliff, N.\ Dominance statistics: ordinal
analyses to answer ordinal questions. \emph{Psychol.\ Bull.} \textbf{114},
494--509 (1993).

\end{thebibliography}

\end{document}
